\chapter{Theoretical Background}
\label{c:literature}

\noindent
This thesis brings together two research strands. One is emotion recognition and the status of
annotator disagreement within it; the other is the framing effect of perceived authorship.
Section~\ref{lit:intersection} shows where the two meet, and Section~\ref{lit:gap} states the gap
this thesis addresses.

\section{Emotion Recognition and Annotator Disagreement}
\label{lit:disagreement}

\noindent
Emotion recognition assigns fine-grained affective labels to text, whereas sentiment analysis
has typically been reduced to polarity. Popular resources such as GoEmotions
\cite{Demszky2020GoEmotions} provide large annotated corpora over rich emotion taxonomies, and transformer classifiers trained on them achieve good
benchmark performance \cite{Acheampong2021EmotionReview}.

Beneath these benchmarks lies a common assumption. Each text is labelled by multiple annotators,
their judgments are aggregated (usually by majority vote), and the aggregate is the gold standard
against which models are trained and scored \cite{Basile2020GoldStandard}. On this view,
annotators are equivalent estimators of one underlying truth, and any difference from the
majority is a mistake.

A large literature argues that this assumption does not hold for subjective tasks.
Pavlick and Kwiatkowski \citeyear{Pavlick2019InherentDisagreement} show that disagreement in
natural language inference is often \emph{inherent} rather than erroneous: annotators differ
genuinely and reproducibly. Uma et al.\ \citeyear{Uma2021Survey} survey learning approaches that
model disagreement instead of eliminating it. Aroyo and Welty \citeyear{Aroyo2015CrowdTruth}
argue directly against the single-ground-truth perspective, and Davani et al.\
\citeyear{Davani2022Disagreement} show that models treating annotators individually can
outperform those trained on aggregated labels.

When disagreement is systematic, aggregation is not neutral. Annotator judgments vary with
demographic background \cite{Prabhakaran2021AnnotatorMetadata}, and in hate speech and
offensive language annotation the variation follows social group membership \cite{Sap2019RacialBias,
Waseem2016AnnotatorInfluence}. Nor is it only demographic. Geva et al.\
\citeyear{Geva2019AnnotatorBias} show that, for datasets whose instances were themselves written
by crowdworkers, a model can learn to identify the most prolific annotators from the text they
wrote, and that the model generalises poorly to examples written by annotators unseen during
training. The identity of the annotator itself, that is, carries systematic, non-interchangeable
signal. A majority label therefore encodes the dominant reading and suppresses minority ones
\cite{Davani2022Disagreement}.

A particularly clear example is emotion. According to Acheampong et al.\
\citeyear{Acheampong2021EmotionReview}, emotional meaning is often implicit and context-dependent,
leaving room for competing but defensible readings. Chou and Lee \citeyear{Chou2019EveryRating}
argue that each rating contains information, so disposing of non-majority ratings would mean
disposing of evidence.

\section{Perceived Authorship and Framing Effects}
\label{lit:framing}

\noindent
The second strand concerns the role that assumptions about a speaker play in interpreting what
is said. The same text can be read differently depending on who is believed to have written
it.

Henestrosa et al.\ \citeyear{Henestrosa2023Authorship} found that, for a science journalism
article held constant in content, authorship cues had an impact on readers' perceptions of the
\emph{author}, and that their experiments did not reveal differences in the article's perceived
credibility or trustworthiness between AI and human authorship. The manipulation is subtle (an
attribution line, not a change in content), but it shifts some judgments and not others. The
present thesis adopts the same design: hold the text fixed, vary only the attributed
author, and quantify the difference, without determining in advance what judgments will vary and
which will not.

Contextual framing sensitivity is well documented for language models. Santurkar et al.\
\citeyear{Santurkar2023WhoseOpinions} demonstrate that model outputs on opinion questions conform
to some demographic groups more than others, and that this shifts when the model is prompted to
adopt a persona. Zhang et al.\ \citeyear{Zhang2018Persona} show that conditioning dialogue
agents on persona profiles systematically changes their responses. Both results establish that
persona information in a prompt is not neutral.

Work applying models to ambiguous emotional content has begun to treat ambiguity as a property in
its own right. Hong et al.\ \citeyear{Hong2024AERLLM} propose ambiguity-aware emotion recognition
with large language models, arguing that a single label is not always an appropriate output
target. R\"ottger et al.\ \citeyear{Rottger2022AnnotationParadigms} distinguish annotation
paradigms aimed at a single, correct answer from those aimed at intentionally representing
variation.

\section{Points of Intersection}
\label{lit:intersection}

\noindent
Neither strand has fully explored the point at which the two meet.

The literature on disagreement shows that variation in subjective annotation is systematic and
worth preserving, but it tends to locate that variation in the \emph{annotator}: their
background, their sensitivity, their interpretive style. The variable is who is reading.

The framing literature shows that assumptions about the \emph{author} shift interpretation.
It usually measures the shift on a single outcome variable such as perceived quality, however,
and typically in humans or in models, not in both under one design.

Together, they raise a question that neither strand asks directly. Suppose disagreement in
emotion annotation is partly driven by differing assumptions about the author, which readers make
silently when context is underspecified. Supplying that context explicitly should then shift the
\emph{distribution} of emotional readings in a measurable and predictable way. If models are
sensitive to persona framing, they should show a similar shift that can be compared against the
human shift on the same texts.

\section{Research Gap and Motivation}
\label{lit:gap}

\noindent
The review above points to three gaps.

\textbf{Perceived authorship is not manipulated as an experimental variable in emotion
recognition.} Existing emotion corpora record disagreement but do not vary author attribution
while holding the text unchanged. Disagreement is observed rather than induced, so its link to
assumptions about the author cannot be isolated.

\textbf{Aggregation destroys the quantity of interest.} If author framing redistributes
probability mass among plausible emotions without necessarily changing which emotion is most
popular, a majority-label analysis is structurally insensitive to it. The effect can only be seen
when judgments are expressed as distributions. This thesis measures directly how much aggregation
loses (Section~\ref{res:vectorization}).

\textbf{Human and model framing effects are rarely measured under one design.} Claims that models
do or do not resemble human interpretation are often based on comparison with a
majority-derived gold label. That tells us whether the model's single answer matches the human's.
It does not tell us whether the model is \emph{responsive to context} as human beings are.
Measuring the change on both sides, under the same manipulation, on the same texts, and comparing
direction and magnitude separately, is a different and narrower question.

\noindent
This thesis takes up all three. It manipulates perceived authorship experimentally, represents
every judgment as a full emotion distribution, and applies the identical manipulation to human
annotators and to seven large language models, so that their responses can be compared directly.
